% T31 — «Алгоритмы и блок-схемы»: каждая задача = блок-схема + условие.
% Информатика, класс 9, 4 задачи на трассировку.
\documentclass[a4paper,11pt]{article}

\usepackage{../../../../Lessons/_templates/preamble_worksheet}
\usepackage{../_shared/worksheet_check}

\usetikzlibrary{shapes.geometric, shapes.misc, arrows.meta, positioning}

% --- Палитра: техно-индиго ---
\definecolor{wcprimary}{HTML}{1F3A93}
\definecolor{wcprimaryLight}{HTML}{B9C5E8}
\definecolor{wcprimaryBG}{HTML}{F0F3FB}
\definecolor{wcprocess}{HTML}{2D7DD2}
\definecolor{wcdecision}{HTML}{F18F01}
\definecolor{wcio}{HTML}{2EB872}

\renewcommand{\familydefault}{\sfdefault}

% Стили узлов блок-схемы
\tikzset{
  bsstart/.style={rounded rectangle, draw=wcprimary, fill=wcprimaryBG,
                  line width=0.7pt, minimum width=18mm, minimum height=6mm, font=\small},
  bsproc/.style ={rectangle, draw=wcprocess, fill=white,
                  line width=0.7pt, minimum width=24mm, minimum height=6mm, font=\small},
  bsdec/.style  ={diamond, draw=wcdecision, fill=white,
                  line width=0.7pt, aspect=2, minimum width=22mm, font=\small},
  bsio/.style   ={trapezium, trapezium left angle=70, trapezium right angle=110,
                  draw=wcio, fill=white, line width=0.7pt, font=\small, minimum width=20mm},
  bsarrow/.style={-stealth, line width=0.8pt, wcprimary!80}}

% --- Заголовок ---
\renewcommand{\WorksheetTitle}[1]{%
  \par\noindent\begin{tikzpicture}
    \fill[wcprimary] (0,0) rectangle (\linewidth, -14mm);
    \node[anchor=west, text=white, font=\bfseries\Large\ttfamily]
      at (3mm, -5mm) {algorithm.trace()};
    \node[anchor=west, text=wcprimaryLight, font=\bfseries\large\sffamily]
      at (3mm, -11mm) {#1};
  \end{tikzpicture}\par\vspace{4mm}}

% TaskBox с блок-схемой
\renewcommand{\TaskBox}[2]{%
  \par\noindent\begin{tikzpicture}
    \node[draw=wcprimary, line width=0.8pt, rounded corners=3pt,
          fill=white, inner sep=3mm,
          text width=\dimexpr\linewidth-6mm\relax,
          anchor=north west] (B) at (0,0) {#2};
    \node[anchor=south west, fill=wcprimary, text=white,
          rounded corners=1pt, inner xsep=5pt, inner ysep=1.5pt,
          xshift=1.5pt] at (B.north west)
      {\scriptsize\bfseries\ttfamily TASK \##1};
  \end{tikzpicture}\par\vspace{4mm}}

\SetAnswerFieldSize{34mm}{9mm}

\begin{document}

\WorksheetTitle{Трассировка алгоритмов}
{\small\itshape\color{wcprimary!85}Информатика \quad·\quad Класс 9 \quad·\quad T31 \quad·\quad ID: T31-algo-2026-05-27}\par\vspace{2mm}
\FirstPageHeader

\TaskBox{1}{%
  \textbf{Найди значение $s$ после выполнения алгоритма.}\par
  Вход: $n = 5$.\par
  \vspace{1.5mm}
  \centering
  \begin{tikzpicture}[node distance=4mm]
    \node[bsstart]                       (n1) {start};
    \node[bsio,    below=of n1]          (n2) {read $n$};
    \node[bsproc,  below=of n2]          (n3) {$s := 0,\ i := 1$};
    \node[bsdec,   below=of n3, yshift=-1mm] (n4) {$i \le n$?};
    \node[bsproc,  right=12mm of n4]     (n5) {$s := s + i;\ i := i + 1$};
    \node[bsio,    below=of n4, yshift=-1mm] (n6) {print $s$};
    \node[bsstart, below=of n6]          (n7) {end};
    \draw[bsarrow] (n1) -- (n2);
    \draw[bsarrow] (n2) -- (n3);
    \draw[bsarrow] (n3) -- (n4);
    \draw[bsarrow] (n4) -- node[above,font=\scriptsize]{да} (n5);
    \draw[bsarrow] (n5) |- ([yshift=1mm]n3.east) -- (n3);
    \draw[bsarrow] (n4) -- node[right,font=\scriptsize]{нет} (n6);
    \draw[bsarrow] (n6) -- (n7);
  \end{tikzpicture}\par
  \vspace{1.5mm}\hfill\textbf{\color{wcprimary}$s=$}~\AnswerField{1}%
}

\TaskBox{2}{%
  \textbf{Какое значение примет переменная $k$?}\par
  Вход: $a = 18,\ b = 12$ (алгоритм Евклида).\par
  \vspace{1.5mm}
  \centering
  \begin{tikzpicture}[node distance=4mm]
    \node[bsstart]                       (m1) {start};
    \node[bsio,   below=of m1]           (m2) {read $a,b$};
    \node[bsdec,  below=of m2, yshift=-1mm] (m3) {$a \ne b$?};
    \node[bsdec,  right=10mm of m3]      (m4) {$a > b$?};
    \node[bsproc, above right=2mm and 8mm of m4] (m5) {$a := a-b$};
    \node[bsproc, below right=2mm and 8mm of m4] (m6) {$b := b-a$};
    \node[bsproc, below=of m3, yshift=-1mm] (m7) {$k := a$};
    \node[bsio,   below=of m7]           (m8) {print $k$};
    \draw[bsarrow] (m1) -- (m2);
    \draw[bsarrow] (m2) -- (m3);
    \draw[bsarrow] (m3) -- node[above,font=\scriptsize]{да}(m4);
    \draw[bsarrow] (m4) -- node[above,font=\scriptsize]{да}(m5);
    \draw[bsarrow] (m4) -- node[below,font=\scriptsize]{нет}(m6);
    \draw[bsarrow] (m5) -| ([xshift=8mm]m3.north) -- (m3);
    \draw[bsarrow] (m6) -| ([xshift=8mm]m3.north) -- (m3);
    \draw[bsarrow] (m3) -- node[right,font=\scriptsize]{нет}(m7);
    \draw[bsarrow] (m7) -- (m8);
  \end{tikzpicture}\par
  \vspace{1.5mm}\hfill\textbf{\color{wcprimary}$k=$}~\AnswerField{2}%
}

\TaskBox{3}{%
  \textbf{Сколько раз выведется слово ``HELLO''?}\par
  $i$ перебирает значения от $1$ до $20$, печать срабатывает только если $i$~--- чётное и $i$ делится на $3$.\par
  \vspace{1.5mm}\hfill\textbf{\color{wcprimary}Ответ:}~\AnswerField{3}%
}

\TaskBox{4}{%
  \textbf{Что будет в массиве $A$ после прохода?}\par
  $A[1..5] = [3, 1, 4, 1, 5]$. Цикл $i=2..5$: $A[i] := A[i] + A[i-1]$.
  Выпиши $A[5]$.\par
  \vspace{1.5mm}\hfill\textbf{\color{wcprimary}$A[5]=$}~\AnswerField{4}%
}

\WorksheetQR{T31-algo/L1/v1}

\end{document}
